\subsection{g-模的对偶}

设 E、F 为 g-模。回顾一下（第 I 章，§3，第 3 节），\(\operatorname{Hom}_{k}(\operatorname{E},\operatorname{F})\) 具有典则的 g-模结构。设 \(\varphi\) 是权为 \(\lambda\) 的 \(\operatorname{Hom}_{k}(\operatorname{E},\operatorname{F})\) 中的元素。若 \(\mu \in h^{*}\)，则 \(\varphi(\operatorname{E}^{\mu}) \subset \operatorname{F}^{\lambda + \mu}\)（第 VII 章，§1，第 1 节，命题 2 (ii)）。因此，若 E 和 F 是有限维的，则权为 \(\lambda\) 的 \(\operatorname{Hom}_{k}(\operatorname{E},\operatorname{F})\) 中的元素，正是《代数》第 II 章 §11 第 2 节定义 4 意义下的次数为 \(\lambda\) 的分次同态。

命题 11. 设 E 是有限维 g-模，并考虑 g-模 \(E^{*} = \operatorname{Hom}_{k}(E, k)\)。

\begin{enumerate}
\def\labelenumi{(\roman{enumi})}
\item
  元素 \(\lambda \in \mathrm{P}\) 是 \(\mathrm{E}^*\) 的权，当且仅当 \(-\lambda\) 是 \(\mathrm{E}\) 的权；并且 \(\lambda\) 在 \(\mathrm{E}^*\) 中的重数等于 \(-\lambda\) 在 \(\mathrm{E}\) 中的重数。
\item
  若 \(\mathrm{E}\) 是最高权为 \(\omega\) 的单模，则 \(\mathrm{E}^*\) 是最高权为 \(-w_0(\omega)\) 的单模（参见第 2 节的备注 2）。
\end{enumerate}

把 k 看作元素权为 0 的平凡 g-模。根据上面的论述，\(E^{*}\) 中权为 \(\lambda\) 的元素，是从 E 到 k 的同态；它们在 \(E^{\mu}\) 上为零，只要 \(\mu \neq -\lambda\)。这证明了 (i)。若 E 是单模，则 \(E^{*}\) 也是单模（第 I 章，§3，第 3 节），最后一个断言由第 2 节的备注 2 得出。

备注。1) 设 E、E* 如命题 11 中，并设 \(\sigma \in \mathrm{Aut}(\mathfrak{g},\mathfrak{h})\) 满足，在 §5 第 1 节的记号下，\(\varepsilon (\sigma) = -w_0\)（§5，第 2 节，命题 2）。设 \(\rho ,\rho^{\prime}\) 是与 E、E* 相关的 \(\mathfrak{g}\) 表示。则 \(\rho \circ \sigma\) 是 \(\mathfrak{g}\) 的一个单表示，其最高权为 \(-w_{0}(\omega)\)，所以 \(\rho \circ \sigma\) 与 \(\rho^{\prime}\) 等价。

\begin{enumerate}
\def\labelenumi{\arabic{enumi})}
\setcounter{enumi}{1}
\tightlist
\item
  假设 \(w_{0} = -1\)。那么对于任意有限维 g-模 E，E 都同构于 \(E^{*}\)。回顾一下，若 g 是单的，则在以下情形中 \(w_{0} = -1\)：g 为 \(A_{1}, B_{l} (l \geq 2)\) 型、\(C_{l} (l \geq 2)\) 型、\(D_{l} (l \text{ even } \geq 4)\) 型，或 \(E_{7}, E_{8}, F_{4}, G_{2}\) 型（第 VI 章，图版）。
\end{enumerate}

引理 2. 令 \(h^0 = \sum_{\alpha \in \mathrm{R}_{+}} H_\alpha\)。则 \(h^0 = \sum_{\alpha \in \mathrm{B}} a_\alpha H_\alpha\)，其中 \(a_\alpha\) 是满足 \(\geq 1\) 的整数。设 \((b_\alpha)_{\alpha \in \mathrm{B}}, (c_\alpha)_{\alpha \in \mathrm{B}}\) 是满足 \(b_\alpha c_\alpha = a_\alpha\) 对所有 \(\alpha \in \mathrm{B}\) 成立的标量族。置 \(x = \sum_{\alpha \in \mathrm{B}} b_\alpha X_\alpha, y = \sum_{\alpha \in \mathrm{B}} c_\alpha X_{-\alpha}\)。存在一个同态 \(\varphi\)，从 \(\mathfrak{sl}(2,k)\) 到 \(\mathfrak{g}\)，使得 \(\varphi(H) = h^0, \varphi(X_+) = x, \varphi(X_-) = y\)。

事实 \(a_{\alpha}\) 是满足 \(\geq 1\) 的整数，是因为 \((H_{\alpha})_{\alpha \in \mathrm{B}}\) 是根系 \((H_{\alpha})_{\alpha \in \mathrm{B}}\) 的一个基（参见第 VI 章，§1，第 5 节，备注 5）。我们有：

\[
\alpha (h ^ {0}) = 2\tag{1}
\]

对所有 \(\alpha \in \mathrm{B}\) 都成立（第 VI 章，§1，第 10 节，命题 29 的推论），因此

\[
[ h ^ {0}, x ] = \sum_ {\alpha \in \mathrm{B}} b _ {\alpha} \alpha (h ^ {0}) X _ {\alpha} = 2 x\tag{2}
\]

\[
[ h ^ {0}, y ] = \sum_ {\alpha \in \mathrm{B}} c _ {\alpha} (- \alpha (h ^ {0})) X _ {- \alpha} = - 2 y.\tag{3}
\]

另一方面，

\[
[ x, y ] = \sum_ {\alpha , \beta \in \mathrm{B}} b _ {\alpha} c _ {\beta} [ X _ {\alpha}, X _ {- \beta} ] = \sum_ {\alpha \in \mathrm{B}} b _ {\alpha} c _ {\alpha} [ X _ {\alpha}, X _ {- \alpha} ] = - \sum_ {\alpha \in \mathrm{B}} a _ {\alpha} H _ {\alpha} = - h ^ {0},\tag{4}
\]

从而同态 \(\varphi\) 存在。

命题 12. 设 E 是有限维单 g-模，\(\omega\) 是其最高权，B 是 E 上 g-不变双线性型的向量空间。设 m 是整数 \(\sum_{\alpha\in\mathrm{R}_{+}}\omega(H_{\alpha})\)，从而 m/2 是 \(\omega\) 关于 B 的坐标之和（第 VI 章，§1，第 10 节，命题 29 的推论）。设 \(w_{0}\) 是满足 \(w_{0}(B) = -B\) 的 W 中的元素。

\begin{enumerate}
\def\labelenumi{(\roman{enumi})}
\item
  若 \(w_0(\omega) \neq -\omega\)，则 \(\mathcal{B} = 0\)。
\item
  假设 \(w_{0}(\omega) = -\omega\)。则 B 的维数为 1，且 B 的每个非零元素都是非退化的。若 m 为偶数（分别为奇数），则 B 的每个元素都是对称的（分别为交错的）。
\end{enumerate}

\begin{enumerate}
\def\labelenumi{\alph{enumi})}
\item
  设 \(\Phi \in \mathcal{B}\)。令 \(\varphi\) 为从 E 到 \(\mathrm{E}^*\) 的映射，对 \(x, y \in \mathrm{E}\) 定义为 \(\varphi(x)(y) = \Phi(x, y)\)；它是 \(\mathfrak{g}\)-模同态。若 \(\Phi \neq 0\)，则 \(\varphi \neq 0\)，所以根据 Schur 引理，\(\varphi\) 是同构，从而 \(\Phi\) 非退化。因此，\(\mathfrak{g}\)-模 E 同构于 \(\mathfrak{g}\)-模 \(\mathrm{E}^*\)，故 \(w_0(\omega) = -\omega\)。至此证明了 (i)。
\item
  从现在起假设 \(w_{0}(\omega) = -\omega\)。则 E 同构于 \(E^{*}\)。向量空间 B 同构于 \(\operatorname{Hom}_{\mathfrak{g}}(\mathrm{E}, \mathrm{E}^{*})\)，从而同构于 \(\operatorname{Hom}_{\mathfrak{g}}(\mathrm{E}, \mathrm{E})\)，后者的维数为 1（§6，第 1 节，命题 1 (iii)）。因此 \(\dim B = 1\)。根据 a)，B 的每个非零元素 \(\Phi\) 都是非退化的。令 \(\Phi_{1}(x, y) = \Phi(y, x)\)，其中 \(x, y \in E\)。由前述结果，存在 \(\lambda \in k\)，使得 \(\Phi_{1}(x, y) = \lambda \Phi(x, y)\) 对所有 \(x, y \in E\) 成立。于是 \(\Phi(y, x) = \lambda \Phi(x, y) = \lambda^{2} \Phi(y, x)\)，所以 \(\lambda^{2} = 1\) 且 \(\lambda = \pm 1\)。因此，\(\Phi\) 或者是对称的，或者是交错的。
\item
根据第 VII 章，§1，第 3 节，命题 9 (v)，\(E^{\lambda}\) 与 \(E^{\mu}\) 关于 \(\Phi\) 正交，只要 \(\lambda + \mu \neq 0\)。由于 \(\Phi\) 非退化，可知 \(E^{\omega}, E^{-\omega}\) 关于 \(\Phi\) 不正交。
\end{enumerate}

\(d\)）存在一个同态 \(\varphi\)，从 \(\mathfrak{sl}(2,k)\) 满射到 \(\mathfrak{g}\) 的某个子代数，并将 \(H\) 映到 \(\sum_{\alpha \in \mathrm{R}_{+}}H_{\alpha}\)（引理 2）。通过这个同态，把 E 看作一个 \(\mathfrak{sl}(2,k)\)-模。于是，\(\mathrm{E}^{\lambda}\) 的元素具有权 \(\lambda \left(\sum_{\alpha \in \mathrm{R}_{+}}H_\alpha\right)\)。若 \(\lambda \in \mathrm{P}\) 满足 \(\mathrm{E}^{\lambda} \neq 0\) 且 \(\lambda \neq \omega, \lambda \neq -\omega\)，则 \(-\omega < \lambda < \omega\)，所以

\[
- m = - \omega \left(\sum_ {\alpha \in \mathrm{R} _ {+}} H _ {\alpha}\right) <   \lambda \left(\sum_ {\alpha \in \mathrm{R} _ {+}} H _ {\alpha}\right) <   \omega \left(\sum_ {\alpha \in \mathrm{R} _ {+}} H _ {\alpha}\right) = m.
\]

令 \(G\) 为类型为 \(V(m)\) 的 \(\mathfrak{sl}(2, k)\)-模 \(E\) 的同构型分支。由上文，\(G\) 的长度为 1，且包含 \(E^{\omega}, E^{-\omega}\)。根据 \(c\)，\(\Phi\) 在 \(G\) 上的限制非零。根据 §1，第 3 项，备注 3，\(m\) 是偶数还是奇数，取决于这个限制是对称的还是交错的。鉴于 \(b\)，证毕。

定义 2. 若 g 的有限维不可约表示 \(\rho\) 在 E 上存在一个关于 \(\rho\) 不变的非退化对称（分别为交错）双线性型，则称该表示是正交型（分别为辛型）的。

